\chapter{方法}\label{chap:method}

本章介绍基于PINN与LLM的统一期权定价框架。在第~\ref{sec:three_models}节中，概述三大定价模型对统一框架设计的关键影响；在第~\ref{sec:unified_pde}节中，介绍统一PDE算子与软掩码机制；在第~\ref{sec:additive_param}节中，介绍加法参数化输出设计；在第~\ref{sec:architecture}节中，描述网络架构；在第~\ref{sec:loss}节中，给出损失函数的构成与权重选择；在第~\ref{sec:training}节中，介绍训练策略；在第~\ref{sec:llm_router}节中，介绍LLM路由层的设计。系统整体流程为：用户以自然语言描述期权需求，LLM路由层解析参数并选择模型，统一PINN求解器返回定价结果。

\section{三大期权定价模型}\label{sec:three_models}

本文涉及的三大期权定价模型已在第~\ref{sec:option_models}节详细介绍，此处仅概述各模型对统一框架设计的关键影响。BSM模型（式~\eqref{eq:bsm}）以常数波动率$\sigma$为唯一随机性来源，存在封闭解析解，本文将其作为加法参数化的基线（见第~\ref{sec:additive_param}节）：
\begin{align}
  V_\text{BS}(S,\tau;\sigma) &= S\Phi(d_1) - Ke^{-r\tau}\Phi(d_2), \label{eq:bs_formula}\\
  d_1 &= \frac{\ln(S/K)+(r+\sigma^2/2)\tau}{\sigma\sqrt{\tau}}, \quad d_2 = d_1 - \sigma\sqrt{\tau}.
\end{align}
CEV模型（式~\eqref{eq:cev}）将扩散系数推广为$\sigma^2 S^{2\beta}$，本文取$\beta=0.5$，无封闭解析解，以Crank-Nicolson有限差分解作为数据锚点。Heston模型（式~\eqref{eq:heston}）引入随机方差$v$，PDE升至二维，以特征函数半解析解作为数据锚点。三个模型的PDE在扩散系数结构上的差异，正是第~\ref{sec:unified_pde}节统一算子设计的出发点。

\section{统一PDE算子}\label{sec:unified_pde}

\subsection{三大模型的PDE统一形式}

BSM、CEV、Heston三大模型的PDE在形式上具有高度相似性，均可写为如下统一算子：
\begin{equation}\label{eq:unified_pde}
  \mathcal{F}[V] =
    \frac{\partial V}{\partial t}
    + \frac{1}{2}a(S,v;\boldsymbol{\lambda})\frac{\partial^2 V}{\partial S^2}
    + b(S,v;\boldsymbol{\lambda})\frac{\partial^2 V}{\partial S\partial v}
    + \frac{1}{2}c(v;\boldsymbol{\lambda})\frac{\partial^2 V}{\partial v^2}
    + rS\frac{\partial V}{\partial S}
    + d(v;\boldsymbol{\lambda})\frac{\partial V}{\partial v}
    - rV = 0,
\end{equation}
其中参数向量$\boldsymbol{\lambda}=(\sigma,\beta,\kappa,\theta,\xi,\rho)$控制各系数，具体对应关系如表~\ref{tab:unified_coefficients}所示。

\begin{table}[htbp]
  \centering
  \caption{统一PDE算子系数与三大模型的对应关系}
  \label{tab:unified_coefficients}
  \begin{tabular}{lcccc}
    \hline
    模型 & $a(S,v;\boldsymbol{\lambda})$ & $b(S,v;\boldsymbol{\lambda})$ & $c(v;\boldsymbol{\lambda})$ & $d(v;\boldsymbol{\lambda})$ \\
    \hline
    BSM & $\sigma^2 S^2$ & $0$ & $0$ & $0$ \\
    CEV & $\sigma^2 S^{2\beta}$ & $0$ & $0$ & $0$ \\
    Heston & $vS^2$ & $\rho\xi vS$ & $\xi^2 v$ & $\kappa(\theta-v)$ \\
    \hline
  \end{tabular}
\end{table}

当$\xi=0$时，$b=c=d=0$，统一算子自动退化为一维BSM/CEV方程；当$\xi>0$时，交叉项和$v$方向扩散项被激活，对应Heston的二维随机波动率结构。

\subsection{软掩码机制}

为实现三大模型之间的连续插值，本文引入软掩码（Soft Mask）：
\begin{equation}\label{eq:soft_mask}
  \text{mask} = \tanh\!\left(\frac{\xi}{0.05}\right)^2.
\end{equation}
当$\xi\to 0$时，$\text{mask}\to 0$；当$\xi$较大（如$\xi=0.3$）时，$\text{mask}\approx 1$。利用软掩码，扩散系数$a$统一为：
\begin{equation}\label{eq:soft_a}
  a = (1-\text{mask})\cdot\sigma^2 S^{2\beta} + \text{mask}\cdot v S^2.
\end{equation}
这一设计使得：当$\xi=0$（BSM/CEV）时，$a=\sigma^2 S^{2\beta}$，精确退化为BSM/CEV扩散项；当$\xi=0.3$（Heston）时，$\text{mask}\approx 0.998$，$a\approx vS^2$，精确对应Heston扩散项；中间状态时，$a$在两者之间连续插值，网络可以学习到平滑的参数依赖关系。软掩码的阈值$0.05$是根据实验调整的超参数：过小会导致BSM/CEV区域的梯度消失，过大会使Heston区域的掩码不够接近1。

\section{加法参数化输出}\label{sec:additive_param}

直接用神经网络输出期权价格存在两个问题：初始化困难（随机初始化的网络输出量级约为$10^{-1}$，而期权价格量级约为$10^0$至$10^2$，训练初期梯度信号弱）；以及深度虚值退化（当$S\ll K$时，乘法参数化$V=V_\text{BS}\cdot(1+\text{net})$在$V_\text{BS}\to 0$时退化，网络无法学习到有效的修正量）。为此，本文采用加法参数化输出：
\begin{equation}\label{eq:additive}
  \hat{V}(S,v,t;\boldsymbol{\lambda}) = V_\text{BS}(S,\tau;\sigma_\text{eff}) + K\cdot\text{net}(x),
\end{equation}
其中$\sigma_\text{eff}$为等效波动率，$\text{net}(x)$为神经网络输出的修正量（以$K$为单位，量级约为$0$至$0.1$）。等效波动率通过软掩码计算：
\begin{equation}\label{eq:sigma_eff}
  \sigma_\text{eff} = (1-\text{mask})\cdot\sigma + \text{mask}\cdot\sqrt{v},
\end{equation}
即BSM/CEV时用常数波动率$\sigma$，Heston时用当前方差$v$的平方根作为近似波动率。网络最后一层权重初始化为$\mathcal{N}(0, 0.001^2)$，使训练开始时$K\cdot\text{net}\approx 0$，网络从接近正确答案的位置开始优化；深度虚值时$V_\text{BS}\to 0$，但$K\cdot\text{net}$仍可输出非零修正，网络可以学习到CEV/Heston相对于BSM的价格差异。

\section{网络架构}\label{sec:architecture}

\subsection{输入特征}

网络输入为9维向量，包含归一化的状态变量和模型参数：
\begin{equation}
  x = \left[\frac{S}{S_\text{max}},\; \frac{v}{v_\text{max}},\; \frac{t}{T},\; \sigma,\; \beta,\; \kappa,\; \theta,\; \xi,\; \rho\right],
\end{equation}
其中$S_\text{max}=300$，$v_\text{max}=1$。归一化使各输入均落在$[0,1]$区间，避免不同量级的输入导致梯度不平衡。

\subsection{网络结构}

网络采用6层全连接结构，每层128个神经元，激活函数为$\tanh$：
\begin{align}
  h_0 &= x \in \mathbb{R}^9, \\
  h_l &= \tanh(W_l h_{l-1} + b_l), \quad l=1,\ldots,6, \\
  \text{net}(x) &= W_7 h_6 + b_7 \in \mathbb{R}.
\end{align}
最后一层权重初始化为$\mathcal{N}(0, 0.001^2)$，偏置初始化为0，确保训练开始时$\text{net}\approx 0$。

$\tanh$激活函数光滑且高阶导数有界，适合PINN中需要计算二阶偏导数的场景，同时其输出范围$(-1,1)$有助于控制网络输出量级。

\subsection{终值条件的近似硬约束}

软约束方式将终值条件纳入损失函数，网络只能近似满足。参考Guo等人\cite{guo2024icpinn}的ICPINN方法，本文通过输出变换将终值条件硬编码。当$t=T$时，$\tau=T-t=0$，Black-Scholes公式自动退化为终值条件$V_\text{BS}(S,0;\sigma_\text{eff})=\max(S-K,0)$；同时，网络最后一层的小初始化确保修正量$K\cdot\text{net}$在训练初期接近0，不破坏终值条件的满足。

\section{损失函数}\label{sec:loss}

统一PINN的总损失由PDE残差、边界条件和数据锚点三部分组成：
\begin{equation}\label{eq:total_loss}
  \mathcal{L}(\theta) = w_\text{pde}\,\mathcal{L}_\text{pde}
                      + w_\text{bc}\,\mathcal{L}_\text{bc}
                      + w_\text{data}\,\mathcal{L}_\text{data}.
\end{equation}

\subsection{PDE残差损失}

在求解域$\Omega\times[0,T)$内随机采样$N_c$个配点$\{(S_i,v_i,t_i)\}$，计算统一PDE算子的残差：
\begin{equation}
  \mathcal{L}_\text{pde} = \frac{1}{N_c}\sum_{i=1}^{N_c}
    \left(\frac{\mathcal{F}[\hat{V}](S_i,v_i,t_i)}{|\hat{V}(S_i,v_i,t_i)|+1}\right)^2.
\end{equation}
分母中的$|\hat{V}|+1$为相对残差归一化项，避免ATM区域（$V\approx 10$）的残差主导损失，使OTM区域（$V\approx 0.1$）也能获得足够的梯度信号。

\subsection{边界条件损失}

边界条件设置为：
\begin{align}
  V(0,v,t) &= 0, \label{eq:bc_lower} \\
  V(S_\text{max},v,t) &= S_\text{max} - Ke^{-r(T-t)}, \label{eq:bc_upper} \\
  V(S,v_\text{max},t) &\approx S. \label{eq:bc_vmax}
\end{align}
式~\eqref{eq:bc_lower}表示股价为0时期权价值为0；式~\eqref{eq:bc_upper}为深度实值近似；式~\eqref{eq:bc_vmax}表示极高波动率下期权价值趋近于标的资产价格。

边界条件损失为：
\begin{equation}
  \mathcal{L}_\text{bc} = \frac{1}{N_b}\sum_{k=1}^{N_b}
    \left|\hat{V}(S_k,v_k,t_k) - g(S_k,v_k,t_k)\right|^2,
\end{equation}
其中$g$为对应边界的目标值。

\subsection{数据驱动锚点损失}

纯PDE约束存在解不唯一的数值问题。本文在$t=0$处预计算参考解（BSM用解析解，CEV用Crank-Nicolson有限差分，Heston用特征函数半解析解），作为数据驱动锚点：
\begin{equation}
  \mathcal{L}_\text{data} = \frac{1}{N_d}\sum_{j=1}^{N_d}
    \left(\hat{V}(S_j,v_j,0) - V_j^*\right)^2,
\end{equation}
其中$V_j^*$为参考解，$N_d=20$（每个模型各20个点）。数据锚点的权重$w_\text{data}=100$远大于PDE权重$w_\text{pde}=1$，确保网络在$t=0$处的精度。

\subsection{损失权重的选择依据}

三个损失项的权重$(w_\text{pde}, w_\text{bc}, w_\text{data})$的选择对训练稳定性和最终精度有重要影响，本文通过以下分析确定权重：

\textbf{量级分析。}PDE残差$\mathcal{L}_\text{pde}$的典型量级为$10^{-2}$至$10^{-1}$（经相对归一化后），边界条件损失$\mathcal{L}_\text{bc}$的典型量级为$10^{-1}$至$10^0$，数据锚点损失$\mathcal{L}_\text{data}$的典型量级为$10^{-2}$至$10^{-1}$。若三者权重相等，数据锚点损失会被边界条件损失主导，导致$t=0$处精度不足。

\textbf{权重设置。}本文设置$w_\text{pde}=1$，$w_\text{bc}=10$，$w_\text{data}=100$。其中$w_\text{data}=100$是因为数据锚点提供了精确的参考解，是最可靠的监督信号，高权重确保网络在已知解处精确拟合，防止PDE约束将解推向错误方向；$w_\text{bc}=10$是因为边界条件是PDE适定性的必要条件，中等权重确保边界处的物理约束被满足，同时不过度压制内部PDE残差；$w_\text{pde}=1$是因为PDE残差是软约束，权重过大会导致训练初期梯度爆炸（因为初始网络的PDE残差很大），权重为1是稳定训练的基准。

\textbf{消融实验验证。}在预实验中，我们测试了$(1,1,1)$、$(1,10,10)$、$(1,10,100)$三组权重，发现$(1,10,100)$在BSM和Heston上均取得最低的ATM相对误差，因此采用此权重组合。

\section{训练策略}\label{sec:training}

\subsection{混合采样}

训练时从三个模型的参数分布中混合采样配点：每个训练步骤中，从BSM、CEV、Heston三个模型各采样$N_c/3$个配点，拼接后计算统一PDE残差。这一策略确保网络在三种模型的参数空间中均匀训练，避免某一模型主导损失。

\subsection{优化器与学习率}

使用Adam优化器\cite{kingma2014adam}，初始学习率$10^{-3}$，余弦退火学习率调度（$\eta_\text{min}=10^{-5}$）。总训练轮数30000轮，在NVIDIA GPU上约需1小时。

\subsection{配点采样策略}

配点在求解域内均匀随机采样：BSM/CEV为一维问题，$S\in[0, S_\text{max}]$，$t\in[0,T]$，$v$固定为$v_0$；Heston为二维问题，$S\in[0, S_\text{max}]$，$v\in[0, v_\text{max}]$，$t\in[0,T]$。每个训练步骤重新采样配点（online sampling），避免过拟合到固定配点集。

\section{LLM路由层}\label{sec:llm_router}

三种定价模型分别适用于不同的市场假设，传统工具要求用户手动指定模型类型和全部数值参数，对非专业用户门槛较高。本文引入DeepSeek LLM作为智能路由层，使用户可以通过自然语言描述期权需求，由LLM自动完成模型选择和参数提取，相比基于规则的解析器，LLM能够处理模糊表达和多语言输入，鲁棒性更强。

系统提示词将LLM角色定义为"期权定价参数提取器"，规定模型选择规则：仅提供$\sigma$则选BSM，提及$\beta$或"CEV"则选CEV，提供$\kappa,\theta,\xi,\rho,v_0$或提及"Heston"/"stochastic volatility"则选Heston，默认BSM。对用户未提及的参数使用美股市场典型默认值（$S=100$，$K=100$，$T=1.0$，$r=0.05$，$\sigma=0.2$，$\beta=0.5$，$v_0=0.04$，$\kappa=2.0$，$\theta=0.04$，$\xi=0.3$，$\rho=-0.7$），并通过\texttt{response\_format: json\_object}强制输出合法JSON，避免自然语言混入结构化输出。LLM的输出遵循以下JSON Schema：
\begin{verbatim}
{
  "model": "bsm" | "cev" | "heston",
  "option_type": "call" | "put",
  "S": float,      // 当前股价
  "K": float,      // 执行价
  "T": float,      // 到期时间（年）
  "r": float,      // 无风险利率
  "sigma": float,  // BSM/CEV波动率参数
  "beta": float,   // CEV弹性参数（仅CEV）
  "v0": float,     // 初始方差（仅Heston）
  "kappa": float,  // 均值回归速度（仅Heston）
  "theta": float,  // 长期方差（仅Heston）
  "xi": float,     // vol-of-vol（仅Heston）
  "rho": float     // 相关系数（仅Heston）
}
\end{verbatim}

系统在接收到JSON输出后，进行参数范围校验（如$\sigma>0$，$\beta\in(0,1]$，$\xi>0$，$\rho\in(-1,1)$），对越界参数截断至合法范围，再传入统一PINN求解器。完整定价流程为：用户输入自然语言描述（中英文均可）；LLM提取参数并选择模型；统一PINN加载预训练权重，构造参数向量$\boldsymbol{\lambda}$；网络前向传播返回期权价格（约2ms）。

\section{本章小结}\label{sec:method_summary}

本章介绍了基于PINN与LLM的统一期权定价框架。核心创新在于：（1）软掩码机制将三大模型的PDE算子统一为单一参数化形式；（2）加法参数化输出以Black-Scholes解析解为基线，解决了深度虚值区域的数值退化问题；（3）近似硬约束机制使终值条件在训练初期即被满足，加速收敛；（4）LLM路由层通过结构化输出实现可靠的参数提取和模型选择。下一章将报告实验结果。
